% !TeX program = lualatex
% !TeX encoding = utf8
% !TeX spellcheck = uk_UA
% !TeX root = ../main.tex


%=========================================================
\Opensolutionfile{answer}[\currfilebase/\currfilebase-Answers]
\Writetofile{answer}{\protect\section*{\nameref*{\currfilebase}}}
\chapter{Основи електростатики діелектриків}\label{\currfilebase}\hypertarget{FieldInDielectrics}{}
\graphicspath{{\currfilebase/Pictures}}
%=========================================================
\epigraph{\Alexander Словом \enquote{діелектрик} я позначаю речовину, крізь яку або через яку діють
	електричні сили.}{Майкл Фарадей, \enquote{Експериментальні дослідження з електрики}, \S\,1168}

%% --------------------------------------------------------
\section{Основні поняття}\hypertarget{dielectrics-basic-concepts}{}
%% --------------------------------------------------------


Досі ми розглядали електричне поле у~вакуумі та у~провідниках. У~провідниках, як ми
з'ясували, вільні заряди~--- електрони~--- під дією як завгодно слабкого поля
переміщуються на макроскопічні відстані, аж доки поле всередині провідника не звернеться
в~нуль. Проте більшість речовин, що оточують нас,~--- скло, вода, повітря, пластмаси,~---
поводяться зовсім інакше: внесені в~електричне поле, вони не гасять його всередині себе,
а~лише \emph{послаблюють}. Такі речовини називають \emphx{діелектриками}.

Відмінність діелектрика від провідника пов'язана з~тим, які заряди є в~його структурі.
Домовимось умовно розділяти всі заряди речовини на два типи.

\begin{enumerate}
	\item \emphx{Вільні заряди}\index{Заряди вільні і зв'язані}~--- це заряди, які можуть переміщатися на великі відстані
	      в~речовині (набагато більші за міжатомні відстані). Саме вільні заряди відповідають
	      за струм і~за повну компенсацію зовнішнього поля всередині провідника.
	      У~діелектриках вільних зарядів, як правило, дуже мало, або немає зовсім.
	\item \emphx{Зв'язані (поляризаційні) заряди}~--- це заряди, які під дією зовнішніх
	      полів або сил лише незначно \emph{зміщуються} відносно свого положення рівноваги
	      (електрон у~атомі, іон у~вузлі кристалічної ґратки) і~повертаються назад,
	      у~положення рівноваги, одразу після зняття зовнішнього впливу.
\end{enumerate}

Отже, \emphx{діелектриками} називають речовини, в~яких є зв'язані заряди, а~вільних зарядів
дуже мало або немає зовсім: усі заряджені частинки міцно утримуються в~межах свого атома чи
молекули.


%% --------------------------------------------------------
\subsection*{Мікроскопічне та макроскопічне поле}
%% --------------------------------------------------------


Щоб коректно описати поведінку речовини в~полі, потрібно спершу з'ясувати, про яке саме
поле йдеться. Якщо подивитись на речовину на атомному масштабі, електричне поле
виявляється надзвичайно нерівномірним: всередині атома чи молекули воно на багато порядків
сильніше, ніж у~проміжку між молекулами, і~різко змінюється від точки до точки та
в~часі~--- адже електрони рухаються, атоми коливаються. Таке істинне, швидкозмінне поле
називають \emphx{мікрополем}\index{Мікрополе і середнє поле} $\Efield_\text{micro}$.

Усі макроскопічні прилади~--- вольтметри, зонди, наші вимірювання~--- фізично не здатні
відстежити настільки швидкі та дрібномасштабні варіації поля. Те, що вони фіксують,~--- це
поле, усереднене по деякому \emph{фізично нескінченно малому} об'єму $\Delta V$: досить
малому порівняно з~розмірами тіла, щоб вважатися \enquote{точкою} в~макроскопічному сенсі,
але водночас достатньо великому, щоб містити величезну кількість атомів і~молекул. Таке
\emphx{середнє (макроскопічне) поле} визначається як
% .........................................................
\begin{equation*}
	\tcbhighmath{
		\left\langle \Efield\right\rangle = \frac1{\Delta V} \iiint\limits_{\Delta V} \Efield_\text{micro}\, dV.
	}
\end{equation*}
% .........................................................

Це поле змінюється в~просторі істотно повільніше, ніж мікрополе, і~саме його ми надалі
й~називатимемо просто $\Efield$, коли говоритимемо про поле всередині речовини.


%% --------------------------------------------------------
\section{Структура діелектриків}\hypertarget{dielectric-structure}{}
%% --------------------------------------------------------


%% --------------------------------------------------------
\subsection*{Дипольні моменти молекул}
%% --------------------------------------------------------


За відсутності зовнішнього поля позитивний заряд ядер атомів, що утворюють молекулу,
і~негативний заряд електронної оболонки в~сумі дають нуль, тобто молекула
електронейтральна. Проте розподіл цих зарядів у~просторі не завжди симетричний. Якщо
\enquote{центр ваги} позитивного заряду молекули не збігається з~\enquote{центром ваги}
негативного, молекула має власний, сталий \emphx{дипольний момент} $\vect p_0$. Причиною
такої асиметрії є \emph{різна електронегативність} атомів, що входять до складу молекули:
атом з~більшою електронегативністю сильніше \enquote{притягує} до себе спільну електронну
хмару, і~на ньому виникає надлишок негативного заряду, тоді як на іншому атомі~---
надлишок позитивного.

Молекули з~ненульовим власним дипольним моментом називають \emphx{полярними}. Класичний
приклад~--- молекула води \ce{H2O}: через кутову форму молекули (кут між зв'язками
\ce{O-H} становить приблизно $105^\circ$) дипольні моменти окремих зв'язків \ce{O-H} не
компенсують одне одного, і~молекула в цілому має досить великий дипольний момент
(рис.~\ref{pic:MoleculeDipoles}).

\begin{figure}[h!]\centering
	\localinput{moleculedipole.tikz}
	\caption{Власні дипольні моменти молекул води та аміаку}
	\label{pic:MoleculeDipoles}
\end{figure}

Дипольні моменти молекул вимірюють у~\emphx{дебаях} (розділ~\ref{ChargeSystem}); одиницю
названо на честь П.~Дебая:
% .........................................................
\begin{equation*}
	1\ \text{Д} = 10^{-18}\ \text{Фр}\cdot\text{см} \approx 3{,}34\cdot10^{-30}\ \text{Кл}\cdot\text{м}.
\end{equation*}
% .........................................................

\begin{table}[h!]
	\centering
	\caption{Дипольні моменти деяких полярних молекул\label{tab:dipoles}}
	\begin{tblr}{
			colspec={X[l, m]X[c,m]},
			row{odd} = {gray!10},
			row{1} = {c, m, fg=white, bg=themecolorlight, font=\bfseries},
			hlines,
			vlines
		}
		Молекула                   & Дипольний момент (Дебай) \\
		Вуглекислий газ (\ce{CO2}) & 0                        \\
		Хлороводень (\ce{HCl})     & 1,08                     \\
		Аміак (\ce{NH3})           & 1,46                     \\
		% Метанол (\ce{CH3OH})       & 1.70                     \\
		Вода (\ce{H2O})            & 1,84                     \\
	\end{tblr}
\end{table}

Більшість полярних молекул має дипольний момент порядку $1$ Д. Це легко зрозуміти:
елементарний заряд, зміщений на типову атомну відстань $1$~Å, дає
$4{,}8\cdot10^{-10}\cdot10^{-8} \approx 4{,}8$~Д, а~в~молекулі зміщується лише частина заряду. Одиниця широко
застосовується у~фізичній хімії, атомній і~молекулярній фізиці. Дипольні моменти деяких
відомих полярних молекул наведено в~таблиці~\ref{tab:dipoles}.

Показовим є приклад вуглекислого газу \ce{CO2}: попри те, що кожен зв'язок \ce{C=O} є
полярним, молекула лінійна і~симетрична, тому дипольні моменти окремих зв'язків повністю
компенсують одне одного, і~сумарний дипольний момент молекули дорівнює нулю.


%% --------------------------------------------------------
\subsection*{Неполярні молекули та наведений дипольний момент}
%% --------------------------------------------------------


\emphx{Неполярні молекули}\index{Молекули полярні і неполярні}~--- це молекули, в~яких електричні заряди розподілені
симетрично, і~власний дипольний момент відсутній ($p_0 = 0$). У~таких молекулах атоми
зазвичай мають однакову або близьку електронегативність, тому спільна електронна хмара
розподілена рівномірно між атомами, або ж молекула має достатньо високу симетрію (як
у~випадку \ce{CO2}). Прикладами неполярних молекул є кисень \ce{O2}, азот \ce{N2},
метан \ce{CH4}, а~також окремі атоми інертних газів.

Важливо розуміти, що відсутність \emphx{власного} дипольного моменту не означає, що
неполярна молекула взагалі не може мати дипольного моменту. У зовнішньому електричному
полі позитивно заряджене ядро і~негативно заряджена електронна хмара зміщуються
в~протилежні боки: ядро~--- за полем, електрони~--- проти поля
(рис.~\ref{pic:InducedDipole}). В~результаті \enquote{центри ваги} позитивного та
негативного зарядів розходяться, і~в~молекулі виникає \emphx{наведений (індукований)
	дипольний момент}, пропорційний прикладеному полю:
% .........................................................
\begin{equation}\label{eq:InducedDipole}
	\vect p = \alpha \Efield,
\end{equation}
% .........................................................
де коефіцієнт $\alpha$ називають \emphx{поляризовністю молекули}.

\begin{figure}[h!]\centering
	\localinput{nonpolar-molecule.tikz}
	\caption{Неполярна молекула без зовнішнього поля та в~зовнішньому
		полі\label{pic:InducedDipole}}
\end{figure}

Цей механізм появи дипольного моменту називають \emphx{електронною (деформаційною)
	поляризацією}. Наведений момент визначається \emph{відносним} зміщенням електронної хмари
та ядра; оскільки електрони дуже легкі, цей механізм встигає за полем навіть на оптичних
частотах. На відміну від
орієнтаційного механізму, який ми розглянемо в~\S~\ref{sec:Langevin-Debye}, електронна
поляризовність практично не залежить від температури.


%% --------------------------------------------------------
\section{Поляризація діелектриків}\label{sec:Polarization}\hypertarget{dielectric-polarization}{}
%% --------------------------------------------------------


\emphx{Поляризацією}\index{Поляризація діелектриків} називають такий перерозподіл зв'язаних
зарядів, за якого речовина набуває дипольного моменту. Найчастіше її спричиняє зовнішнє
електричне поле, але можуть і~інші чинники~--- механічна деформація, зміна температури
(\S~\ref{sec:PiezoPyro}).

Розглянемо, як відбувається поляризація в~двох основних класах діелектриків.

\begin{itemize}
	\item \emphx{Полярні діелектрики}\index{Діелектрики!полярні} складаються з~полярних молекул, кожна з~яких уже має
	      власний дипольний момент $\vect p_0 \neq 0$. За відсутності зовнішнього поля теплове
	      (хаотичне) обертання молекул орієнтує їхні дипольні моменти цілком випадково в~усіх
	      напрямках, тому в~будь-якому фізично малому об'ємі дипольні моменти в~середньому
	      взаємно компенсуються, і~сумарний дипольний момент речовини дорівнює нулю. Коли ж
	      прикладене зовнішнє поле, воно створює момент сил, що прагне розвернути кожен диполь
	      уздовж поля (див.~формулу $\vect M = \vect p\times\Efield$, отриману
	      в~розділі~\ref{ChargeSystem}). Тепловий рух заважає ідеальному впорядкуванню, проте певна
	      переважна орієнтація вздовж поля все ж встановлюється~--- виникає так звана
	      \emphx{орієнтаційна поляризація}\index{Поляризація діелектриків!орієнтаційна} (рис.~\ref{pic:PolarizationPolar}).

	      % ==================== FIGURE 1: полярні молекули ====================
	      \begin{figure}[h!]
		      \centering
		      % --- (а) кадр 1: полярні молекули без поля ---
		      \begin{subfigure}[t]{0.48\linewidth}
			      \centering
			      \localinput{PolarNoField.tikz}
			      \caption{Без поля}
			      \label{pic:PolarNoField}
		      \end{subfigure}\hfill
		      % --- (б) кадр 2: полярні молекули в полі ---
		      \begin{subfigure}[t]{0.48\linewidth}
			      \centering
			      \localinput{PolarField.tikz}
			      \caption{У~полі}
			      \label{pic:PolarField}
		      \end{subfigure}
		      \caption{Поляризація полярних молекул:
		      (\protect\subref{pic:PolarNoField})~--- без поля;
		      (\protect\subref{pic:PolarField})~--- у~зовнішньому електричному полі,
		      диполі орієнтуються вздовж поля.}
		      \label{pic:PolarizationPolar}
	      \end{figure}

	\item \emphx{Неполярні діелектрики}\index{Діелектрики!неполярні} складаються з~неполярних молекул, у~яких за
	      відсутності поля центри позитивного й~негативного зарядів збігаються. Зовнішнє поле
	      зміщує електронні хмари відносно ядер~\eqref{eq:InducedDipole}, і~в~кожній
	      молекулі виникає наведений дипольний момент, зорієнтований уздовж поля, ---
	      \emphx{електронна поляризація}\index{Поляризація діелектриків!електронна} (рис.~\ref{pic:PolarizationNonPolar}).

	      % ==================== FIGURE 2: неполярні молекули ====================
	      \begin{figure}[h!]
		      \centering
		      % --- (а) кадр 3: неполярні молекули без поля ---
		      \begin{subfigure}[t]{0.48\linewidth}
			      \centering
			      \localinput{NonPolarNoField.tikz}
			      \caption{Без поля}
			      \label{pic:NonPolarNoField}
		      \end{subfigure}\hfill
		      % --- (б) кадр 4: неполярні молекули в полі ---
		      \begin{subfigure}[t]{0.48\linewidth}
			      \centering
			      \localinput{NonPolarField.tikz}
			      \caption{У~полі}
			      \label{pic:NonPolarField}
		      \end{subfigure}
		      \caption{Поляризація неполярних молекул:
		      (\protect\subref{pic:NonPolarNoField})~--- без поля;
		      (\protect\subref{pic:NonPolarField})~--- у~зовнішньому електричному полі
		      молекули деформуються, утворюючи індукований дипольний момент.}
		      \label{pic:PolarizationNonPolar}
	      \end{figure}
\end{itemize}

У~випадку іонних кристалів (наприклад, \ce{NaCl}) існує ще один механізм~---
\emphx{іонна поляризація}\index{Поляризація діелектриків!іонна}: позитивні та негативні іони кристалічної ґратки під дією поля
зміщуються один відносно одного як цілі підґратки, що також створює об'ємний дипольний
момент. Детальніше про всі три механізми (електронний, іонний, орієнтаційний) і~про те, як
кожен із них залежить від частоти поля й~температури, йтиметься
в~\S~\ref{sec:AtomPolarizability}.

В~обох випадках~--- полярних і~неполярних діелектриків~--- результат один: на межах
виділеного в~речовині об'єму з'являються надлишкові (некомпенсовані) зв'язані заряди, тоді
як в~основному об'ємі речовини позитивні та негативні заряди сусідніх диполів взаємно
компенсують одне одного. Саме ці надлишкові поверхневі заряди і~створюють власне поле
речовини $\Efield'$, яке, за принципом суперпозиції, додається до зовнішнього поля
$\Efield_\text{ex}$:
% .........................................................
\begin{equation*}
	\Efield = \Efield_\text{ex} + \Efield'.
\end{equation*}
% .........................................................
Оскільки диполі орієнтуються так, що їхні позитивні полюси \enquote{дивляться} за полем,
а~негативні~--- проти поля, власне поле $\Efield'$ усередині діелектрика, як правило,
спрямоване \emph{назустріч} зовнішньому, і~сумарне поле \emph{послаблюється}. Наскільки
саме~--- залежить від форми тіла та його орієнтації. У~пластині, перпендикулярній до поля,
поле слабшає в~$\epsilon$ разів, у~кулі~--- в~$(\epsilon+2)/3$ раза, а~в~пластині,
паралельній полю, не змінюється взагалі: зв'язані заряди виступають лише на її вузьких
торцях і~всередині пластини поля майже не створюють (задача~\ref{prb:DielectricSlab}).

На рисунку~\ref{pic:PolarizedDielectric} схематично показано поляризований діелектрик
у~зовнішньому електричному полі: на протилежних поверхнях виділеного об'єму виникають
некомпенсовані зв'язані заряди $-Q$ та $+Q$, власне поле яких спрямоване назустріч
зовнішньому.

\begin{figure}[h!]
	\centering
	\localinput{PolarizedDielectric.tikz}
	\caption{Поляризований діелектрик у~зовнішньому електричному полі: на поверхнях
		виникають зв'язані заряди $-Q$ та $+Q$. Стрілки означають:
		\protect\tikz[baseline=-0.6ex]{\protect\draw[->, red, thick, opacity=0.5] (0,0) --
			(0.6,0);} --- зовнішнє поле $\Efield_\text{ex}$;
		\protect\tikz[baseline=-0.6ex]{\protect\draw[<-, red, thick, dashed] (0,0) --
			(0.6,0);} --- поле зв'язаних зарядів $\Efield'$;
		\protect\tikz[baseline=-0.6ex]{\protect\draw[->, brown, thick] (0,0) -- (0.6,0);} ---
		вектор поляризації $\vect{P}$ (збігається за напрямком із~результуючим полем $\Efield$). }
	\label{pic:PolarizedDielectric}
\end{figure}


%% --------------------------------------------------------
\subsection*{Електрострикція}
%% --------------------------------------------------------


Уміщення діелектрика в~електричне поле призводить не лише до поляризації, а~й~до
\emph{зміни розмірів тіла}~--- явища, яке називають \emphx{електрострикцією}\index{Електрострикція}.
Електрострикція спостерігається в~усіх без винятку діелектриках, проте в~переважній
більшості матеріалів виражена слабко: деформація пропорційна \emph{квадрату} поля,
$\Delta \ell \sim E^2$, тому не залежить від напрямку поля (зміна знака поля не змінює
знака деформації) і стає помітною лише в~дуже сильних полях. У~техніці електрострикцію
(у~кристалах без центра симетрії) \enquote{перекриває} значно сильніший і~лінійний за
полем \emphx{п'єзоефект}\index{П'єзоефект} (\S~\ref{sec:PiezoPyro}), тому як самостійне
явище електрострикція застосовується рідше.


%% --------------------------------------------------------
\section{Вектор поляризації}\hypertarget{polarization-vector}{}
%% --------------------------------------------------------


Для кількісного опису стану поляризованого діелектрика вводять \emph{макроскопічну}
характеристику~--- \emphx{вектор поляризації}\index{Вектор поляризації} (або \emphx{поляризованість}), який дорівнює
дипольному моменту одиниці об'єму речовини:
% .........................................................
\begin{equation}\label{eq:Pdef}
	\tcbhighmath{
		\vect P = \frac1V \sum_i \vect p_i,
	}
\end{equation}
% .........................................................
де підсумовування проводиться за всіма молекулярними диполями $\vect p_i$, що потрапляють
у~фізично малий об'єм $V$.

Для великого класу ізотропних діелектриків (гази, більшість рідин та аморфних твердих тіл)
у~не надто сильних полях експериментально встановлено, що поляризованість лінійно залежить від напруженості поля
\emph{в~самому діелектрику}:
% .........................................................
\begin{equation}\label{eq:PlinearE}
	\tcbhighmath{
		\vect P = \chi \Efield,
	}
\end{equation}
% .........................................................
де безрозмірна величина $\chi > 0$ називається \emphx{діелектричною
	сприйнятливістю}\index{Діелектрична сприйнятливість} речовини. Сприйнятливість~---
макроскопічний аналог \emph{поляризовності} окремої молекули $\alpha$
(формула~\eqref{eq:InducedDipole}): $\alpha$ описує відгук однієї молекули, $\chi$~---
одиниці об'єму речовини, а~нетривіальний зв'язок між ними~--- предмет
\S~\ref{sec:ClausiusMossotti}. (Термін \enquote{поляризовність} ми надалі вживаємо лише
для атомів і~молекул.) Сприйнятливість $\chi$ не залежить від $\Efield$ і~характеризує властивості
самого діелектрика (тип молекул, їхню концентрацію, температуру тощо). Такі діелектрики
називають \emph{лінійними}.

Не всі діелектрики лінійні. В~анізотропних кристалах сприйнятливість є \emphx{тензором}:
$P_i = \sum_j \chi_{ij}E_j$, і~поле вздовж однієї осі кристала поляризує його й~уздовж
інших. У~дуже сильних (наприклад, лазерних) полях до лінійного доданка додаються квадратичні
й~вищі, $P_i = \sum_j \chi^{(1)}_{ij}E_j + \sum_{jk}\chi^{(2)}_{ijk}E_jE_k + \ldots$~--- на
цьому ґрунтується нелінійна оптика. У~\emphx{сегнетоелектриках} зв'язок $\vect P$ з~$\Efield$
нелінійний і~неоднозначний (гістерезис), а~в~\emph{електретах} поляризація існує й~без поля
(\S~\ref{sec:FE}).


%% --------------------------------------------------------
\subsection*{Однорідна та неоднорідна поляризація}
%% --------------------------------------------------------


Поляризацію називають \emphx{однорідною}, якщо вектор $\vect P$ є сталим у~всьому об'ємі
речовини, $\vect P = \const$, і~\emphx{неоднорідною}, якщо $\vect P$ змінюється від точки
до точки.

\emphx{Однорідна поляризація.} Нехай поляризація однорідна. Розглянемо косокутний
паралелепіпед, вирізаний із поляризованої речовини так, що його бічне ребро паралельне
вектору поляризації $\vect P$ (рис.~\ref{pic:PolarizedParallelepiped}). Якщо $S$~--- площа
основи (торцевої грані), а~$\ell$~--- довжина бічного ребра паралелепіпеда, то його об'єм
дорівнює $V = S\ell\cos\theta$, де $\theta$~--- кут між ребром $\vect \ell$ і~нормаллю до
основи.

\begin{SCfigure}[1][h!]\centering
	\localinput{parallelepiped.tikz}
	\caption{Косокутний паралелепіпед, вирізаний з~однорідно поляризованого діелектрика
		уздовж вектора поляризації $\vect P$\label{pic:PolarizedParallelepiped}}
\end{SCfigure}

Диполі всередині такого паралелепіпеда шикуються \enquote{голова до хвоста}: позитивний
заряд одного диполя сусідить з~негативним зарядом наступного, і~в об'ємі речовини ці
заряди взаємно компенсуються. Некомпенсованими лишаються тільки заряди на двох торцевих
гранях: на одній виступає надлишковий позитивний заряд, на іншій~--- негативний. Тому весь
паралелепіпед еквівалентний одному великому диполю з~дипольним моментом $\vect p = \sigma'
	S \vect \ell$, де $\sigma'$~--- поверхнева густина \emphx{зв'язаного (поляризаційного)}
заряду на торцевій грані. З~іншого боку, за означенням~\eqref{eq:Pdef}, дипольний момент
паралелепіпеда дорівнює $\vect p = \vect P \cdot V = \vect P \, S\ell\cos\theta$.
Прирівнюючи обидва вирази і~враховуючи, що $\vect P \cdot \vect n = P\cos\theta$ (де
$\vect n$~--- зовнішня нормаль до грані), отримуємо
% .........................................................
\begin{equation}\label{eq:SigmaPrime}
	\tcbhighmath{
		\vect P \cdot \vect n = P_n = \sigma'.
	}
\end{equation}
% .........................................................

Нормальна складова вектора поляризації на поверхні діелектрика дорівнює поверхневій
густині зв'язаного заряду. Там, де $\vect P$ спрямований \emph{із} речовини, на поверхні
виступає додатний зв'язаний заряд, а~там, де $\vect P$ спрямований \emph{у} речовину,~---
від'ємний.

\emphx{Неоднорідна поляризація.} Нехай тепер поляризація неоднорідна, і~в~товщі речовини
виділено об'єм $V$ довільної форми, обмежений замкненою поверхнею $S$. Кожен малий елемент
поверхні $dS$ можна, як і~вище, вважати гранню елементарного косокутного паралелепіпеда,
тому на ньому виступає поверхневий зв'язаний заряд $d q' = \sigma' dS = P_n\, dS$
(рис.~\ref{pic:NonUniformPolarization}).

\begin{SCfigure}[1][htbp]\centering
	\localinput{nonuniform-polarization.tikz}
	\caption{Неоднорідно поляризований діелектрик: вектор $\vect P$ (коричневі стрілки)
		зростає від центру. Крізь поверхню $S$ виходить більше зв'язаного заряду, ніж
		входить, тому на ній виступає $\sigma' = P_n > 0$, а~в~об'ємі $V$ лишається
		від'ємний зв'язаний заряд $\rho' = -\divg\vect P$\label{pic:NonUniformPolarization}}
\end{SCfigure}

Підсумовуючи (інтегруючи) цей вираз за всією замкненою поверхнею, знаходимо повний заряд,
що вийшов на зовнішню поверхню виділеного об'єму:
% .........................................................
\begin{equation*}
	q'_{\text{пов}} = \oiint\limits_S P_n\, dS.
\end{equation*}
% .........................................................
Але сама речовина, з~якої вирізано цей об'єм, електронейтральна: поляризація лише
\emph{перерозподіляє} наявні заряди, не створюючи нових. Тому якщо на зовнішню поверхню
вийшов заряд $q'_{\text{пов}}$, то всередині об'єму мусив залишитися зв'язаний заряд
протилежного знака, розподілений з об'ємною густиною $\rho'$:
% .........................................................
\begin{equation}\label{eq:GaussForP-int}
	\tcbhighmath{
		\oiint\limits_S \vect P\cdot d\vect S = \oiint\limits_S P_n\, dS = - \iiint\limits_V \rho'\, dV.
	}
\end{equation}
% .........................................................
Це співвідношення називають \emphx{теоремою Гаусса для вектора поляризації}\index{Гаусса
теорема!для вектора поляризації}: воно за формою повністю аналогічне теоремі Гаусса для
поля $\Efield$ (\S~\ref{sec:Gauss}), проте зв'язує потік вектора $\vect P$ не з~повним,
а~лише зі \emphx{зв'язаним} зарядом, узятим із протилежним знаком.

Так само як ми переходили від інтегральної форми теореми Гаусса для $\Efield$ до
диференціальної, стягуючи об'єм $V$ у~точку, з формули~\eqref{eq:GaussForP-int} отримуємо
\emphx{диференціальну форму теореми Гаусса для вектора поляризації}:
% .........................................................
\begin{equation}\label{eq:GaussForP-diff}
	\tcbhighmath{
		\divg \vect P = -\rho'.
	}
\end{equation}
% .........................................................
Дивергенція вектора поляризації в~даній точці дорівнює об'ємній густині зв'язаного заряду
в~цій точці, узятій зі знаком мінус. У~окремому випадку однорідної поляризації, коли
$\vect P = \const$, маємо $\divg \vect P \equiv 0$, а~отже, і~$\rho' = 0$: об'ємний
зв'язаний заряд в~однорідно поляризованому діелектрику відсутній, увесь надлишковий заряд
зосереджений на поверхні, що узгоджується з~розглянутим вище прикладом паралелепіпеда.


%% --------------------------------------------------------
\section{Вектор електричної індукції}\label{sec:Dfield}\hypertarget{electric-displacement-vector}{}
%% --------------------------------------------------------


Електричне\index{Вектор електричної індукції} поле всередині речовини створюється не лише
вільними, а~й~зв'язаними (поляризаційними) зарядами. Тому в~теоремі Гаусса для поля
$\Efield$ (яка, нагадаємо, є універсальним законом і~виконується для будь-яких зарядів)
слід урахувати обидва типи зарядів~--- і~вільні $\rho$, і зв'язані $\rho'$:
% .........................................................
\begin{equation*}
	\divg \Efield = 4\pi(\rho + \rho').
\end{equation*}
% .........................................................
Скориставшись знайденим щойно співвідношенням~\eqref{eq:GaussForP-diff}, а~саме $\rho' =
	-\divg\vect P$, отримуємо:
% .........................................................
\begin{equation*}
	\divg \Efield = 4\pi\rho - 4\pi \divg\vect P
	\quad\Longrightarrow\quad
	\divg\left(\Efield + 4\pi \vect P\right) = 4\pi\rho.
\end{equation*}
% .........................................................
Вираз у~дужках зручно позначити окремим символом і~назвати \emphx{вектором електричної
	індукції}:
% .........................................................
\begin{equation}\label{eq:Ddef}
	\tcbhighmath{
		\Dfield = \Efield + 4\pi\vect P.
	}
\end{equation}
% .........................................................
Тоді теорема Гаусса для електричного поля \emph{в~речовині} набуває вигляду:
% .........................................................
\begin{equation}\label{eq:GaussForD-diff}
	\tcbhighmath{
		\divg\Dfield = 4\pi\rho,
	}
\end{equation}
% .........................................................
а~в~інтегральній формі~--- відповідно
% .........................................................
\begin{equation}\label{eq:GaussForD-int}
	\tcbhighmath{
		\oiint\limits_S \Dfield\cdot d\vect S = 4\pi \iiint\limits_V \rho\, dV.
	}
\end{equation}
% .........................................................

\begin{Attention}
	У~теорему Гаусса для $\Dfield$ входять \emph{тільки вільні заряди}: внесок зв'язаних
	уже \enquote{захований} в~означенні $\Dfield$. Тому для потоку $\Dfield$ досить знати
	розподіл вільних зарядів, які зазвичай задані (наприклад, заряд обкладок конденсатора).
\end{Attention}

Підкреслимо: на відміну від $\Efield$, вектор $\Dfield$ \emph{не} є силовим полем~--- він
не діє безпосередньо на пробний заряд. Це допоміжна, розрахункова величина. Те, що
\emph{джерелами} $\Dfield$ є лише вільні заряди, треба розуміти саме в~сенсі дивергенції:
від вільних зарядів залежить потік $\Dfield$ крізь замкнену поверхню. Саме ж поле
$\Dfield$ залежить і~від форми діелектрика: оскільки $\rot\Dfield = 4\pi\rot\vect P$, на
межах діелектрика $\Dfield$, узагалі кажучи, має \enquote{вихори} і~не є потенціальним.
Тому знайти $\Dfield$ лише за вільними зарядами можна тільки в~задачах із високою
симетрією (сферичною, циліндричною, плоскою), коли для цього досить теореми Гаусса.
Приклад протилежної ситуації~--- електрет (\S~\ref{sec:FE}): вільних зарядів немає зовсім,
а~поле $\Dfield$ є.


%% --------------------------------------------------------
\subsection*{Діелектрична проникність}
%% --------------------------------------------------------


Для лінійних діелектриків, у~яких справедливе співвідношення~\eqref{eq:PlinearE}, вектор
індукції теж виявляється пропорційним $\Efield$:
% .........................................................
\begin{equation*}
	\Dfield = \Efield + 4\pi\vect P = \Efield + 4\pi\chi\Efield = (1+4\pi\chi)\Efield.
\end{equation*}
% .........................................................
Позначаючи
% .........................................................
\begin{equation}\label{eq:eps}
	\tcbhighmath{
		\epsilon = 1 + 4\pi\chi,
	}
\end{equation}
% .........................................................
отримуємо просте й~дуже часто використовуване співвідношення
% .........................................................
\begin{equation}\label{eq:DepsE}
	\tcbhighmath{
		\Dfield = \epsilon\Efield.
	}
\end{equation}
% .........................................................
Величину $\epsilon$ називають \emphx{діелектричною проникністю}\index{Діелектрична
проникність} середовища. Для звичайних (не сегнетоелектричних) речовин у~постійних полях
$\chi > 0$, тому $\epsilon > 1$: діелектрична проникність будь-якого
\enquote{нормального} діелектрика перевищує одиницю.
Для вакууму, в~якому поляризація відсутня ($\chi = 0$), маємо $\epsilon = 1$, і~вектори
$\Dfield$ та $\Efield$ збігаються.

Для порівняння наведемо порядки величин $\epsilon$ для кількох типових речовин (постійне
поле, кімнатна температура):
\begin{center}
	\begin{tblr}{
			colspec={X[l, m]X[c,m]},
			row{odd} = {gray!10},
			row{1} = {c, m, fg=white, bg=themecolorlight, font=\bfseries},
			hlines,
			vlines
		}
		Речовина                      & $\epsilon$     \\
		Повітря (сухе)                & $1{,}0006$     \\
		Вода                          & $81$           \\
		Скло                          & $5$--$7$       \\
		Титанат барію (\ce{BaTiO3}) & $10^3$--$10^4$
	\end{tblr}
\end{center}
Бачимо, що для газів $\epsilon$ дуже близька до одиниці (молекули далеко одна від одної),
для рідин і~твердих тіл~--- суттєво більша, а~для сегнетоелектриків (\S~\ref{sec:FE})
сягає аномально великих значень.

\begin{figure}[tbp]\centering
	\begin{subfigure}[t]{0.48\linewidth}\centering
		\resizebox{\linewidth}{!}{\localinput{sphere_Efield.tikz}}
		\caption{Поле $\Efield$}
		\label{pic:SphereE}
	\end{subfigure}\hfill
	\begin{subfigure}[t]{0.48\linewidth}\centering
		\resizebox{\linewidth}{!}{\localinput{sphere_Dfield.tikz}}
		\caption{Поле $\Dfield$}
		\label{pic:SphereD}
	\end{subfigure}
	\caption{Діелектрична куля ($\epsilon = 5$) в~однорідному полі: силові лінії
		$\Efield$~(\protect\subref{pic:SphereE}) і~$\Dfield$~(\protect\subref{pic:SphereD}).
		Усередині кулі обидва поля однорідні; лінії $\Efield$ рідшають (поле слабшає), лінії
		$\Dfield$ згущуються й~неперервно проходять крізь поверхню, бо вільних зарядів на ній
		немає. Знаками показано зв'язані заряди на поверхні}
	\label{pic:DielectricSphere}
\end{figure}

\begin{Example}{Діелектрична куля в~однорідному полі}
	Перевіримо твердження з~\S~\ref{sec:Polarization}, що в~кулі поле слабшає в~$(\epsilon+2)/3$
	раза (рис.~\ref{pic:DielectricSphere}). Спершу знайдемо поле \emph{однорідно поляризованої} кулі радіуса $R$. Таку кулю можна
	уявити як дві кулі з~об'ємними зарядами $+\rho$ і~$-\rho$, зміщені одна відносно одної на
	малий вектор $\vect d$, так що $\vect P = \rho\vect d$: у~спільній частині заряди
	компенсуються, а~на поверхні лишається зв'язаний заряд $\sigma' = P_n$. Усередині
	рівномірно зарядженої кулі поле $\frac{4\pi}{3}\rho\,\vect r$~\eqref{eq:BallField}, тож для
	двох куль із~центрами в~$\pm\vect d/2$
	\begin{equation*}
		\Efield' = \frac{4\pi}{3}\rho\left(\vect r - \frac{\vect d}{2}\right)
		- \frac{4\pi}{3}\rho\left(\vect r + \frac{\vect d}{2}\right)
		= -\frac{4\pi}{3}\vect P.
	\end{equation*}
	Поле зв'язаних зарядів усередині кулі \emph{однорідне} і~напрямлене проти $\vect P$.

	Тепер нехай діелектрична куля вміщена в~однорідне зовнішнє поле $\Efield_0$. Припустимо,
	що вона поляризується однорідно. Тоді поле всередині $\Efield = \Efield_0 -
	\frac{4\pi}{3}\vect P$ теж однорідне, і~з~$\vect P = \chi\Efield$
	\begin{equation*}
		\Efield = \frac{\Efield_0}{1 + 4\pi\chi/3} = \frac{3}{\epsilon+2}\,\Efield_0 .
	\end{equation*}
	Припущення узгоджене: однорідне поле справді дає однорідну поляризацію, і~за теоремою
	єдиності це і~є розв'язок. Поле в~кулі слабшає в~$(\epsilon+2)/3$ раза~--- менше, ніж
	у~пластині, перпендикулярній до поля ($\epsilon$ разів), бо зв'язані заряди кулі
	розподілені по її поверхні, а~не зосереджені на двох гранях. Зовні ж поляризована куля
	створює поле диполя з~моментом $\frac{4\pi}{3}R^3\vect P$ (задача~\ref{prb:PolarBall}).
	Порівняйте з~металевою кулею (розділ~\ref{FieldInMetalls}): там $\epsilon\to\infty$, поле
	всередині нульове, а~дипольний момент $R^3\Efield_0$.
\end{Example}

\emph{Приклад.} Розглянемо точковий заряд $Q$, вміщений у~нескінченний однорідний
діелектрик з~проникністю $\epsilon$ (рис.~\ref{pic:PointChargeDielectric}). Через сферичну
симетрію задачі, за теоремою Гаусса для вектора $\Dfield$, як і~у~випадку поля у~вакуумі,
% .........................................................
\begin{equation*}
	\Dfield = \frac{Q}{r^2}\frac{\vect r}{r},
\end{equation*}
% .........................................................
а~звідси, користуючись~\eqref{eq:DepsE}, знаходимо напруженість поля безпосередньо
в~діелектрику:
% .........................................................
\begin{equation*}
	\Efield = \frac{\Dfield}{\epsilon} = \frac{Q}{\epsilon r^2}\frac{\vect r}{r}.
\end{equation*}
% .........................................................
Оскільки $\epsilon > 1$, поле точкового заряду в~діелектрику всюди \emph{слабкіше}, ніж
таке саме поле у~вакуумі. Наприклад, у~воді ($\epsilon \approx 81$) поле на тій самій
відстані від заряду слабкіше у~$81$ раз, а~у~склі ($\epsilon \approx 5$)~--- у~$5$ разів.
Це прямий кількісний прояв уже відомого нам якісного результату: \emph{поляризаційні
	заряди, що виникають навколо заряду $Q$, послаблюють його поле}. Справді, біля самого
заряду $Q$ зосереджується зв'язаний заряд протилежного знака $-Q(\epsilon-1)/\epsilon$
(задача~\ref{prb:BoundAroundQ}), тож здалеку заряд \enquote{виглядає} як $Q/\epsilon$.

\begin{SCfigure}[2][h!]\centering
	\localinput{point-charge-dielectric.tikz}
	\caption{Точковий заряд у~нескінченному однорідному
		діелектрику\label{pic:PointChargeDielectric}}
\end{SCfigure}


%% --------------------------------------------------------
\section{Граничні умови в~присутності діелектриків}\hypertarget{dielectric-boundary-conditions}{}
%% --------------------------------------------------------


Як\index{Граничні умови!на межі діелектриків} і~у~випадку поля у~вакуумі
(\S~\ref{sec:BC_vacuum}), на межі розділу двох середовищ (двох діелектриків, або
діелектрика й~вакууму) поле може змінюватись стрибком. Знайдемо, як пов'язані значення
полів по обидва боки межі.


%% --------------------------------------------------------
\subsection*{Перша гранична умова (для нормальної складової $\Dfield$)}
%% --------------------------------------------------------


Застосуємо теорему Гаусса для вектора $\Dfield$ до нескінченно малого циліндра, що охоплює
невелику ділянку межі розділу середовищ 1 та 2, так само, як ми робили це для поля
$\Efield$ на поверхні провідника (рис.~\ref{pic:BC1}). Нехай $\sigma$ --- поверхнева
густина \emphx{вільного} заряду на межі розділу (якщо він там є). Потоком через бічну
поверхню циліндра, висота якого прямує до нуля, нехтуємо. Враховуючи, що $d\vect
	S_2 = -d\vect S_1 = \vect n\, dS$, де $\vect n$ --- нормаль, спрямована з~середовища 1
у~середовище 2, маємо

\begin{SCfigure}[1][h!]\centering
	\localinput{bc1.tikz}
	\caption{Циліндр Гаусса на межі розділу двох середовищ для виведення граничної умови
		на нормальну складову $\Dfield$\label{pic:BC1}}
\end{SCfigure}
% .........................................................
\begin{equation*}
	\oiint\limits_S \Dfield\cdot d\vect S = 4\pi\sigma \, dS
	\quad\Longrightarrow\quad
	\Dfield_2 \cdot d\vect S_2 + \Dfield_1\cdot d\vect S_1 \equiv (D_{2n} - D_{1n})\, dS = 4\pi\sigma\, dS,
\end{equation*}
% .........................................................
звідки
% .........................................................
\begin{equation}\label{eq:BC1-D}
	\tcbhighmath{
		D_{2n} - D_{1n} = 4\pi\sigma.
	}
\end{equation}
% .........................................................

\begin{Attention}
	Без вільних зарядів на межі ($\sigma = 0$) нормальна складова $\Dfield$ неперервна,
	$D_{1n} = D_{2n}$, а~нормальна складова $\Efield$ однаково змінюється стрибком~--- через
	зв'язані поверхневі заряди.
\end{Attention}


%% --------------------------------------------------------
\subsection*{Друга гранична умова (для тангенціальної складової $\Efield$)}
%% --------------------------------------------------------


Ця умова не пов'язана з~наявністю діелектрика і~виводиться так само, як і раніше
(\S~\ref{sec:BC_vacuum}), із теореми про циркуляцію вектора $\Efield$ уздовж нескінченно
малого прямокутного контуру $L$, що проходить по обидва боки межі розділу
(рис.~\ref{pic:BC2}); ця теорема залишається справедливою в~будь-якому середовищі~--- адже
електростатичне поле завжди потенціальне, незалежно від того, вакуум це чи діелектрик:

\begin{SCfigure}[1][h!]\centering
	\localinput{bc2.tikz}
	\caption{Контур циркуляції на межі розділу двох середовищ для виведення граничної
		умови на тангенціальну складову $\Efield$\label{pic:BC2}}
\end{SCfigure}
% .........................................................
\begin{equation}\label{eq:BC2-E}
	\tcbhighmath{
		E_{1\tau} = E_{2\tau}.
	}
\end{equation}
% .........................................................
Тангенціальна складова $\Efield$ однакова по обидва боки межі розділу і~не зазнає стрибка.


%% --------------------------------------------------------
\subsection*{Заломлення силових ліній на границі діелектрика}
%% --------------------------------------------------------


Розглянемо межу розділу двох однорідних лінійних діелектриків із проникностями
$\epsilon_1$ та $\epsilon_2$, на якій вільні заряди відсутні. Нехай силова лінія поля
$\Efield_1$ (а~отже, і~$\Dfield_1$) у~середовищі 1 утворює кут $\alpha_1$ з~нормаллю до
межі, а~в~середовищі 2~--- кут $\alpha_2$ (рис.~\ref{pic:RefractionE}). Граничні
умови~\eqref{eq:BC1-D} (без вільних зарядів) та~\eqref{eq:BC2-E} дають:

\begin{figure}[h!]\centering
	\localinput{refraction.tikz}
	\caption{Заломлення силових ліній $\Efield$ і~$\Dfield$ на межі розділу двох
		діелектриків\label{pic:RefractionE}}
\end{figure}
% .........................................................
\begin{equation*}
	\begin{cases}
		E_{1\tau} = E_{2\tau} \quad\Longrightarrow\quad E_1\sin\alpha_1 = E_2\sin\alpha_2, \\
		D_{1n} = D_{2n} \quad\Longrightarrow\quad \epsilon_1 E_1\cos\alpha_1 = \epsilon_2 E_2\cos\alpha_2.
	\end{cases}
\end{equation*}
% .........................................................
Поділивши перше рівняння на друге, отримуємо \emphx{закон заломлення силових ліній}\index{Заломлення силових ліній}:
% .........................................................
\begin{equation}\label{eq:RefractionE}
	\tcbhighmath{
		\frac{\tg\alpha_2}{\tg\alpha_1} = \frac{\epsilon_2}{\epsilon_1}.
	}
\end{equation}
% .........................................................

\begin{Attention}
	У~середовищі з~\emph{більшою} $\epsilon$ лінії поля утворюють \emph{більший} кут
	із~нормаллю, тобто \enquote{притискаються} до межі. Це схоже на заломлення світла, але
	закон зв'язує тангенси, а~не синуси, і~діє в~протилежний бік: в~оптично густішому
	середовищі промінь, навпаки, наближається до нормалі.
\end{Attention}


%%% --------------------------------------------------------
\section{Мікроскопічна теорія діелектричної проникності}\hypertarget{microscopic-permittivity-theory}{}
%%% --------------------------------------------------------


Ми ввели діелектричну проникність $\epsilon$ як \emph{феноменологічну} величину: ми
знаємо, як її виміряти (через ємність конденсатора) і~як нею користуватися в~розрахунках.
З'ясуємо тепер, як $\epsilon$ пов'язана з~властивостями окремих молекул.


%% --------------------------------------------------------
\subsection*{Механізми поляризовності молекули}\label{sec:AtomPolarizability}
%% --------------------------------------------------------


У~зовнішньому полі $\Efield$ молекула набуває наведеного дипольного моменту $\vect p =
	\alpha\Efield$ (формула~\eqref{eq:InducedDipole}). Але \emph{звідки} береться цей
момент~--- питання не одне, а~фізично різне залежно від будови молекули:

\begin{itemize}
	\item $\alpha_e$~--- \emphx{електронна}: поле зміщує легку електронну хмару відносно
	      ядра. Є \emph{в~усіх} атомах і~молекулах без винятку, встигає за полем аж до
	      оптичних частот ($\sim10^{15}$ Гц);
	\item $\alpha_i$~--- \emphx{іонна}: у~іонних кристалах поле зміщує одні іони відносно
	      інших (маса набагато більша за електронну) --- \enquote{вимикається} вже
	      в~ІЧ-діапазоні;
	\item $\alpha_\text{ор}$~--- \emphx{орієнтаційна}: якщо молекула \emph{вже має}
	      власний дипольний момент $p_0$ (полярна молекула), поле не наводить,
	      а~\emph{розвертає} його. Найповільніший і~найчутливіший до $T$ механізм~---
	      \enquote{вимикається} вже у~НВЧ-діапазоні.
\end{itemize}

Повна поляризовність на даній частоті~--- сума тих внесків, які встигають за полем.
У~статичному полі діють усі три: $\alpha = \alpha_e + \alpha_i + \alpha_\text{ор}$.


%% --------------------------------------------------------
\subsection*{Орієнтаційна поляризація полярних діелектриків}\label{sec:Langevin-Debye}
%% --------------------------------------------------------


Полярна молекула ($p_0\neq0$) в~полі $\Efield$ має потенціальну енергію
% .........................................................
\begin{equation*}
	U(\theta) = -\vect p_0\cdot\Efield = -p_0 E\cos\theta
\end{equation*}
% .........................................................
(\S~\ref{sec:DipoleEnergy}), де $\theta$~--- кут між диполем і~полем; енергія мінімальна
($U=-p_0E$) при орієнтації вздовж поля, максимальна ($U=+p_0E$)~--- проти поля. За
відсутності поля тепловий рух орієнтує диполі цілком хаотично; у~полі ж число молекул,
орієнтованих в~одиничний тілесний кут навколо напрямку $\theta$, підкоряється
\emphx{розподілу Больцмана} ($n_0$~--- стала нормування, $n_0 = n/4\pi$ за відсутності поля):
% .........................................................
\begin{equation}\label{eq:Boltzmann}
	n(\theta) = n_0\, e^{-U/kT} = n_0\, e^{+p_0 E\cos\theta / kT}.
\end{equation}
% .........................................................
Це й~означає порушення хаосу: молекул, орієнтованих вздовж поля ($\cos\theta=1$, менша
енергія), \emph{експоненційно більше}, ніж орієнтованих проти поля ($\cos\theta=-1$).

У~практично важливому випадку~--- не надто сильне поле й~не надто низька температура,
$p_0E \ll kT$~--- показник експоненти малий, тож розподіл~\eqref{eq:Boltzmann} можна
розкласти в~ряд і~обмежитись лінійним членом:
% .........................................................
\begin{equation}\label{eq:BoltzmannLinear}
	n(\theta) \approx n_0\left(1 + \frac{p_0 E \cos\theta}{kT}\right).
\end{equation}
% .........................................................
Фізично: за кімнатної температури $kT \approx 4\cdot10^{-14}$~ерг, а~для $p_0 = 1$~Д
у~полі $E = 10^5$~В/см ($\approx 330$ од. СГС) енергія орієнтації $p_0E \approx
3\cdot10^{-16}$~ерг~--- приблизно на два порядки менша. Тому переважна орієнтація~--- \emph{слабка}
поправка до майже хаотичного розподілу, а~не повне вишикування диполів вздовж поля.

Поляризація $P$~--- це дипольний момент одиниці об'єму. У~напрямку, перпендикулярному
полю, внески різних диполів у~середньому компенсуються (за симетрією
$\theta\to\pi-\theta$); вздовж поля~--- ні. Тому $P$ напрямлена вздовж поля і~дорівнює
добутку концентрації молекул $n$ на середню проєкцію дипольного моменту:
% .........................................................
\begin{equation*}
	P = n\, p_0\, \overline{\cos\theta} = n p_0 \cdot \frac{\displaystyle\int_0^\pi \cos\theta \, n(\theta)\, 2\pi\sin\theta\, d\theta}
	{\displaystyle\int_0^\pi n(\theta)\, 2\pi\sin\theta\, d\theta}.
\end{equation*}
% .........................................................
Підставляючи сюди лінеаризований розподіл~\eqref{eq:BoltzmannLinear} і~виконуючи
інтегрування за кутом (у~знаменнику лишається $n_0$, у чисельнику~--- лише доданок
з~$\cos\theta$, оскільки $\int\cos\theta\,d\Omega=0$, а~$\int\cos^2\theta\,d\Omega =
	4\pi/3$), одержуємо
% .........................................................
\begin{equation}\label{eq:LangevinDebye}
	\tcbhighmath{
		P = \frac{n p_0^2}{3kT}\, E.
	}
\end{equation}
% .........................................................
Це і~є \emphx{закон Ланжевена--Дебая}\index{Ланжевена--Дебая закон}: поляризація речовини лінійна за полем, а~коефіцієнт
пропорційності явно виражений через мікроскопічні характеристики молекули ($p_0$)
і~речовини ($n$, $T$).

Оскільки за означенням $P = \chi E$, порівняння з~\eqref{eq:LangevinDebye} одразу дає
сприйнятливість речовини (а~заразом і~орієнтаційну поляризовність \emph{однієї} молекули,
$\chi = n\alpha_\text{ор}$):
% .........................................................
\begin{equation*}
	\chi = \frac{n p_0^2}{3kT}
	\qquad\Longrightarrow\qquad
	\alpha_\text{ор} = \frac{p_0^2}{3kT}.
\end{equation*}
% .........................................................
Підставляючи $\chi$ в~означення проникності $\epsilon = 1+4\pi\chi$, знаходимо
макроскопічний результат:
% .........................................................
\begin{equation}\label{eq:CurieLaw}
	\tcbhighmath{
		\epsilon = 1 + \frac{4\pi n p_0^2}{3kT}.
	}
\end{equation}
% .........................................................
Тут ми врахували лише орієнтаційний внесок і~знехтували електронним: для сильно полярних
молекул $\alpha_e$ значно менше за $p_0^2/3kT$ (задача~\ref{prb:WaterVapour}). З~його
урахуванням $\epsilon = 1 + 4\pi n\left(\alpha_e + p_0^2/3kT\right)$. Обидві формули
справедливі для газів, де поле, що діє на молекулу, практично збігається із середнім
(\S~\ref{sec:ClausiusMossotti}).

Залежність $\alpha_\text{ор} \sim 1/T$ називають \emphx{законом Кюрі}\index{Кюрі закон}: що вища
температура, то сильніше тепловий рух руйнує впорядкування диполів. Тому помітне спадання
$\epsilon$ з~ростом $T$~--- пряма ознака того, що діелектрик полярний: електронна й~іонна
поляризовність від температури практично не залежать.


%% --------------------------------------------------------
\subsection*{Неполярні діелектрики. Формула Клаузіуса--Моссотті}\label{sec:ClausiusMossotti}
%% --------------------------------------------------------


Для розріджених середовищ справджується просте співвідношення $\chi = n\alpha$: кожна
молекула відчуває лише зовнішнє (середнє) поле $\Efield$, бо наведений момент окремої
молекули за означенням $\alpha$ (формула~\eqref{eq:InducedDipole}) визначається саме тим
полем, яке \emph{реально діє} на неї, $\vect p = \alpha\Efield_\text{лок}$, а~не середнім
полем $\Efield$ у~речовині. У~розрідженому газі відстані між молекулами великі, тож поле,
яке створюють на місці обраної молекули всі \emph{інші}, вже поляризовані нею
молекули,~--- нехтовно мале порівняно із зовнішнім: $\Efield_\text{лок}\approx\Efield$.
Тому поляризація речовини (дипольний момент одиниці об'єму з~концентрацією $n$)~--- це
просто
% .........................................................
\begin{equation*}
	\vect P = n\vect p = n\alpha\Efield_\text{лок} \approx n\alpha\Efield,
\end{equation*}
% .........................................................
звідки, порівнюючи з~означенням $\vect P = \chi\Efield$, і~виходить $\chi \approx
	n\alpha$.

У~густому ж середовищі кожна молекула занурена в~поле \emph{вже поляризованих} сусідів,
тож реальне (локальне) поле $\Efield_\text{лок}$, що на неї діє,~--- не $\Efield$, а~деяка
відмінна від нього величина, яку й~належить знайти.

Оточимо обрану молекулу уявною сферою Лоренца (яка мала порівняно з~тілом, однак велика
порівняно з~міжмолекулярними відстанями). Поле, що діє на молекулу,~--- це сума трьох внесків
(рис.~\ref{pic:CavityField}): середнє поле $\Efield$; мінус поле кулі-вирізу,
$+\frac{4\pi}{3}\vect P$ (поле однорідно поляризованої кулі з~протилежним знаком, див.
приклад у~\S~\ref{sec:Dfield} і~задачу~\ref{prb:hole_in_dielectric}); і~поле дискретних
сусідів усередині сфери. Для кубічних кристалів, газів і~слабополярних рідин останній
внесок за симетрією дорівнює нулю, звідки

% .........................................................
\begin{equation}\label{eq:LorentzField}
	\tcbhighmath{
		\Efield_\text{лок} = \Efield + \frac{4\pi}{3}\vect P.
	}
\end{equation}
% .........................................................

Це наближення добре працює для кубічних кристалів і~більшості неполярних та слабополярних
речовин. Для сильно полярних рідин з~водневими зв'язками (вода) його треба уточнювати
(теорія Онзагера).

\begin{SCfigure}[2][h!]\centering
	\localinput{CavityField.tikz}
	\caption{Сферична порожнина Лоренца в~поляризованому
		діелектрику\label{pic:CavityField}}
\end{SCfigure}

Маючи локальне поле~\eqref{eq:LorentzField}, запишемо наведений момент молекули саме через
нього: $\vect p = \alpha\Efield_\text{лок}$. Тоді для поляризації (дипольного моменту
одиниці об'єму з~концентрацією молекул $n$) маємо
% .........................................................
\begin{equation*}
	\vect P = n\vect p = n\alpha\Efield_\text{лок} = n\alpha\left(\Efield + \frac{4\pi}{3}\vect P\right).
\end{equation*}
% .........................................................
Це~--- рівняння відносно самої $\vect P$: вона входить і~в~ліву, і~в~праву частину, бо
локальне поле саме поляризацією середовища й~визначається (задача самоузгоджена~---
поляризація створює поле, яке, своєю чергою, впливає на поляризацію). Розкриваючи дужки та
групуючи доданки з~$\vect P$:
% .........................................................
\begin{equation*}
	\vect P\left(1 - \frac{4\pi n\alpha}{3}\right) = n\alpha\,\Efield
	\quad\Longrightarrow\quad
	\vect P = \frac{n\alpha}{1 - \dfrac{4\pi n\alpha}{3}}\,\Efield.
\end{equation*}
% .........................................................
Порівнюючи це з~означенням сприйнятливості речовини $\vect P = \chi\Efield$, знаходимо
мікроскопічний вираз для $\chi$, що вже враховує локальне поле:
% .........................................................
\begin{equation}\label{eq:ChiFromAlpha}
	\chi = \frac{n\alpha}{1 - \dfrac{4\pi n\alpha}{3}}.
\end{equation}
% .........................................................

Знайдемо діелектричну проникність $\epsilon = 1+4\pi\chi$. З~\eqref{eq:ChiFromAlpha}
% .........................................................
\begin{equation*}
	\epsilon - 1 = 4\pi\chi = \frac{4\pi n\alpha}{1 - \dfrac{4\pi n\alpha}{3}}, \qquad
	\epsilon + 2 = (\epsilon - 1) + 3 = \frac{3}{1 - \dfrac{4\pi n\alpha}{3}}.
\end{equation*}
% .........................................................
Поділивши перший вираз на другий, отримаємо:
% .........................................................
\begin{equation}\label{eq:ClausiusMossotti}
	\tcbhighmath{
		\frac{\epsilon - 1}{\epsilon + 2} = \frac{4\pi}{3}\, n\alpha.
	}
\end{equation}
% .........................................................
Цей вираз називається \emphx{формулою Клаузіуса--Моссотті}\index{Клаузіуса--Моссотті
формула}. Її головна цінність~--- вона \emph{напряму} зв'язує макроскопічно вимірювану
величину ($\epsilon$) з мікроскопічною молекулярною характеристикою ($\alpha$)
і~концентрацією $n$, причому враховує (на відміну від наївної $\chi = n\alpha$) ефекти
взаємної поляризації сусідніх молекул. Використовують її в~обидва боки: за відомим
$\alpha$ обчислюють $\epsilon$, або, навпаки, за виміряним $\epsilon$ й~відомою
концентрацією $n$ знаходять $\alpha$~--- один з~небагатьох прямих макроскопічних методів
вимірювання мікроскопічної характеристики речовини.

Для розрідженого газу ($\frac{4\pi}{3}n\alpha\ll1$) знаменник у~\eqref{eq:ChiFromAlpha}
можна замінити одиницею, і~формула зводиться до $\chi\approx n\alpha$, з~якого ми
й~почали; поправка на локальне поле стає суттєвою вже в~густій рідині чи твердому тілі.

Оцінимо, коли поправка суттєва. Поляризовність $\alpha \sim 10^{-24}$~см$^3$ (порядку
атомного об'єму). Для газу ($n\approx2{,}7\cdot10^{19}$~см$^{-3}$)
$\frac{4\pi}{3}n\alpha\sim10^{-4}$~--- нехтовно мало, а~для рідини ($n\sim3\cdot10^{22}$~см$^{-3}$)
вже $\sim10\%$.


%% ========================================================
\section{Типи діелектриків та їх застосування}\hypertarget{dielectric-types-applications}{}
%% ========================================================


%% --------------------------------------------------------
\subsection*{Електрети}
%% --------------------------------------------------------


\emphx{Електрети}\index{Електрети}~--- це діелектрики, здатні тривалий час зберігати
поляризацію \emph{після зняття} зовнішнього поля, яке цю поляризацію створило. Електрет
часто називають \enquote{\emph{електричним аналогом постійного магніту}}: подібно до того,
як постійний магніт створює магнітне поле без зовнішнього джерела струму, заряджений
електрет створює електричне поле $\Efield$ і~поле індукції $\Dfield$ поза своєю поверхнею
навіть за відсутності будь-якого зовнішнього поля (рис.~\ref{pic:Electret}).

\begin{figure}[h!]\centering
	\localinput{electret-fields.tikz}
	\caption{Поля зарядженого електрета\label{pic:Electret}}
\end{figure}

В~ідеальному електреті вільних зарядів немає, тож $\divg\Dfield = 0$ усюди. Це не означає,
що поля $\Dfield$ немає: так само $\divg\Bfield \equiv 0$, а~магнітне поле існує. Джерелом
полів $\Efield$ і~$\Dfield$ є зв'язані заряди електрета, а~\enquote{вихори} $\Dfield$,
$\rot\Dfield = 4\pi\rot\vect P$, зосереджені на його поверхні.

Один з~поширених способів виготовлення електретів полягає в~тому, що діелектрик нагрівають
і~одночасно поміщають у~сильне електричне поле, так щоб полярні молекули встигли
переорієнтуватись і~вишикуватись уздовж поля. Якщо потім діелектрик охолодити (не знімаючи
поля), теплова рухливість молекул різко падає, повернення диполів у~хаотичний стан
ускладнюється, і~після зняття зовнішнього поля поляризація зберігається протягом місяців,
років, а~іноді й~десятиліть. Такий стан називають \enquote{\emphx{замороженою}
	поляризацією}.

Перші електрети виготовив японський фізик М.~Егучі (1919) із суміші карнаубського воску
й~каніфолі. У~сучасній техніці
більшість електретів виготовляють із синтетичних полімерів --- фторопластів,
поліпропілену, поліетилентерефталату (ПЕТ) тощо. Електрети широко застосовуються
в~електретних мікрофонах, дозиметрах, датчиках тиску та інших пристроях, де потрібне
джерело сталого електричного поля без витрат енергії на його підтримку.


%% --------------------------------------------------------
\subsection*{Сегнетоелектрики}\label{sec:FE}
%% --------------------------------------------------------


\emphx{Сегнетоелектрик}\index{Сегнетоелектрики}~--- це кристалічний діелектрик, який має
\emphx{спонтанну поляризацію всередині окремих доменів}, що виникає \emph{без} будь-якого
зовнішнього електричного поля. \emphx{Домен}\index{Домени!сегнетоелектричні}~--- це
локальна область усередині сегнетоелектрика (типового розміру від кількох до сотень
мікрометрів), у~межах якої дипольні моменти всіх елементарних комірок кристала мимовільно
вишикувані в~одному напрямку, тобто поляризація в~домені стала (рис.~\ref{pic:DomainsE}).
Однак у~різних доменах напрямки поляризації, як правило, різні, тому сумарна поляризація
незбудженого зразка в~цілому близька до нуля.

\begin{SCfigure}[1.5][h!]\centering
	\includegraphics[width=0.4\linewidth]{domains}
	\caption{Доменна структура сегнетоелектрика\label{pic:DomainsE}}
\end{SCfigure}

Під впливом зовнішнього поля межі доменів (доменні стінки) зміщуються, а~самі домени
переорієнтовуються так, щоб їхня поляризація узгоджувалася з~напрямком зовнішнього поля;
домени, поляризація яких збігається з~полем, ростуть за рахунок сусідніх доменів. У~результаті
сумарна поляризація зразка зростає значно швидше й~до значно більших значень, ніж
у~звичайних лінійних діелектриках.

Ця перебудова доменної структури має важливу особливість: вона \emph{неоднозначна}
і~зберігає \enquote{пам'ять} про попередню історію зразка. Залежність поляризації $P$ від
поля $E$ при циклічній зміні останнього утворює характерну петлю, яку називають
\emph{петлею гістерезису} (рис.~\ref{pic:HysteresisE}):

\begin{figure}[tbp]\centering
	\localinput{hysteresis.tikz}
	\caption{Петля гістерезису сегнетоелектрика: $P_r$~--- залишкова поляризація, $E_c$
		--- коерцитивне поле\label{pic:HysteresisE}}
\end{figure}
\begin{enumerate}
	\item за досить сильного поля $E$ поляризація виходить на насичення~--- усі домени вже
	      переорієнтовані, і~подальше зростання $P$ з~полем відбувається практично лінійно, як
	      у~звичайного діелектрика;
	\item якщо тепер зменшувати поле до нуля, поляризація \emph{не} повертається в~нуль,
	      а~залишається на певному рівні $P_r$, який називають \emphx{залишковою
		      поляризацією};
	\item щоб звести поляризацію до нуля, потрібно прикласти поле протилежного знака~---
	      $-E_c$; величину $E_c$ називають \emphx{коерцитивним полем}\index{Коерцитивне поле}
	      (за аналогією з~магнетиками~--- також коерцитивною силою);
	\item при подальшому зростанні поля протилежного знака поляризація знову виходить на
	      насичення у~протилежний бік, а~при поверненні поля до нуля залишається на рівні
	      $-P_r$,~--- і~цикл замикається.
\end{enumerate}
Ці властивості спостерігаються лише в~певному діапазоні температур, нижчих за характерну
для кожного сегнетоелектрика \emphx{температуру Кюрі}\index{Кюрі температура}: вище цієї
температури тепловий рух руйнує спонтанну доменну структуру, і~речовина переходить
у~звичайний (параелектричний) стан.

Найвідоміші сегнетоелектрики~--- титанат барію \ce{BaTiO3} і~сегнетова сіль
\ce{KNaC4H4O6*4H2O} (саме на честь останньої й~названо весь цей клас речовин). Завдяки
аномально великим значенням діелектричної проникності $\epsilon$ (яка для
сегнетоелектриків може сягати $\sim 10^3$--$10^4$) сегнетоелектрики застосовують
у~мініатюрних конденсаторах великої ємності та варикондах (нелінійних конденсаторах зі
змінною ємністю), а~завдяки сильному зв'язку між поляризацією й~механічною деформацією~---
у~широкому спектрі електромеханічних і~механоелектричних перетворювачів: датчиках
мікропереміщень, гідрофонах, акселерометрах, стабілізаторах частоти тощо.


%% --------------------------------------------------------
\subsection*{Інші типи діелектриків}\label{sec:PiezoPyro}
%% --------------------------------------------------------


Окрім електретів і~сегнетоелектриків, до \emphx{активних} діелектриків~--- тобто таких,
чиї властивості (спонтанна поляризація, гістерезис, зв'язок поляризації з~механічною дією
чи температурою) самі по собі є функціональними~--- відносять ще кілька класів речовин.
Розглянемо їх коротко.

\emphx{П'єзоелектрики}\index{П'єзоелектрики}~--- кристали без центра симетрії, в~яких
деформація й~поляризація пов'язані лінійно (прямий і~зворотний ефекти); приклад~--- кварц;
застосування~--- резонатори, п'єзодатчики, ультразвукові перетворювачі.
\emphx{Піроелектрики}\index{Піроелектрики}~--- кристали зі спонтанною поляризацією, що
змінюється з~температурою (звідси заряд на поверхні); усі піроелектрики~--- п'єзоелектрики
(не навпаки), а~всі сегнетоелектрики~--- піроелектрики; застосовуються в~ІЧ-детекторах
і~тепловізорах.

\emphx{Антисегнетоелектрики} (напр., \ce{PbZrO3}) мають антипаралельну поляризацію
сусідніх комірок (сумарно~--- нуль), але полем стрибком переводяться у~сегнетоелектричний
стан без залишку~--- застосовуються в~імпульсних накопичувачах енергії.

\emphx{Рідкі кристали}\index{Рідкі кристали}~--- органічні речовини зі стрижнеподібними
молекулами, що утворюють мезоморфну фазу: плинні, як рідина, орієнтаційно впорядковані, як
кристал; анізотропна сприйнятливість ($\chi_\parallel\neq\chi_\perp$) дає змогу полю
переорієнтовувати молекули~--- на цьому ґрунтується робота LCD-дисплеїв.

Натомість у~переважній більшості технічних застосувань діелектрик відіграє \emph{пасивну}
роль ізолятора або середовища конденсатора, і~його властивості цінують не за нелінійністю
чи керованістю, а~за високою електричною міцністю та малими втратами. До поширених
пасивних діелектриків належать: тверді~--- слюда, кераміка, фторопласт, поліпропілен;
рідкі~--- трансформаторне мастило; газоподібні~--- елегаз \ce{SF6}.


%% --------------------------------------------------------
\section*{Підсумки}
\phantomsection\addcontentsline{toc}{section}{Підсумки}
%% --------------------------------------------------------

На відміну від провідників, у~діелектриках майже немає вільних зарядів: усі заряди
\emphx{зв'язані} й~у~полі лише незначно зміщуються. Такий перерозподіл зв'язаних зарядів
називають \emphx{поляризацією}. Вона має три мікроскопічні механізми: \emphx{електронний}
(зміщення електронної хмари, наведений момент $\vect p = \alpha\Efield$, є в~усіх
речовинах), \emphx{іонний} (зміщення підґраток у~іонних кристалах) та
\emphx{орієнтаційний} (розворот власних дипольних моментів $\vect p_0$ полярних молекул
проти теплового хаосу). Надлишкові зв'язані заряди створюють власне поле $\Efield'$, яке
всередині діелектрика, як правило, напрямлене назустріч зовнішньому, тому поле
в~діелектрику \emph{послаблюється} (наскільки~--- залежить від форми тіла).
Макроскопічним полем в~речовині називають поле, усереднене по фізично нескінченно малому
об'єму.

Кількісно поляризацію описує \emphx{вектор поляризації} $\vect P$~--- дипольний момент
одиниці об'єму; для лінійних діелектриків $\vect P = \chi\Efield$, де $\chi$~---
\emphx{діелектрична сприйнятливість} (макроскопічний аналог поляризовності молекули
$\alpha$). Нормальна складова $\vect P$
на поверхні дорівнює густині зв'язаного заряду, $P_n = \sigma'$, а~в~об'ємі $\divg\vect P
	= -\rho'$. Введення вектора індукції $\Dfield = \Efield + 4\pi\vect P$ дає змогу записати
теорему Гаусса в~речовині через \emphx{лише вільні заряди}: $\divg\Dfield = 4\pi\rho$,
а~для лінійного діелектрика $\Dfield = \epsilon\Efield$, де $\epsilon = 1 + 4\pi\chi > 1$.
Зокрема, поле точкового заряду в~однорідному діелектрику слабкіше, ніж у~вакуумі,
в~$\epsilon$ разів. На межі розділу двох середовищ виконуються граничні умови $D_{2n} -
	D_{1n} = 4\pi\sigma$ (де $\sigma$~--- густина вільного заряду) та $E_{1\tau} = E_{2\tau}$,
звідки випливає закон заломлення силових ліній $\tg\alpha_2/\tg\alpha_1 =
	\epsilon_2/\epsilon_1$.

Мікроскопічна теорія пов'язує $\epsilon$ з~властивостями молекул. Для полярних
діелектриків розподіл Больцмана дає закон Ланжевена--Дебая $P = np_0^2E/(3kT)$ і, як
наслідок, для полярних газів $\epsilon \approx 1 + 4\pi np_0^2/(3kT)$: проникність спадає
з~температурою (закон Кюрі). У~густому середовищі на молекулу діє не середнє, а~\emphx{локальне} поле
$\Efield_\text{лок} = \Efield + \frac{4\pi}{3}\vect P$, і~замість наївного $\chi =
	n\alpha$ справджується формула Клаузіуса--Моссотті $(\epsilon - 1)/(\epsilon + 2) = 4\pi
	n\alpha/3$. Окремі класи діелектриків мають особливі властивості: \emphx{електрети}
зберігають поляризацію після зняття поля, \emphx{сегнетоелектрики} мають доменну
структуру, гістерезис і~аномально велике $\epsilon$ нижче температури Кюрі, а~п'єзо- та
піроелектрики пов'язують поляризацію з~деформацією й~температурою. Ми навчилися описувати
поле в~речовині; наступний розділ присвячено енергії електростатичного поля
(розділ~\ref{ElEnergy}).


%% --------------------------------------------------------
\section*{Задачі}
\phantomsection \addcontentsline{toc}{section}{Задачі}
%% --------------------------------------------------------

% задачі цього розділу короткі: не розриваємо їх між сторінками

\begin{problem}\label{prb:PolarBall}
	Куля радіуса $R$ має однорідну поляризацію $\vect{P}$. Знайдіть густину зв'язаного заряду
	на її поверхні та поле зовні кулі. Покажіть, що зовні воно збігається з~полем точкового
	диполя в~центрі. (Скористайтеся моделлю двох зміщених заряджених куль з~прикладу
	в~\S~\ref{sec:Dfield}.)
	\begin{solution}
		$\sigma' = P\cos\theta$; зовні дві зміщені кулі діють як точкові заряди $\pm Q = \pm\frac43\pi R^3\rho$ на відстані $d$, тобто як диполь $\vect p = Q\vect d = \frac{4\pi}{3}R^3\vect P$ у~центрі.
	\end{solution}
\end{problem}

\begin{problem}\label{prb:hole_in_dielectric}
	У~діелектричному середовищі з~діелектричною проникністю $\epsilon$ існує однорідне поле
	з~напруженістю $\Efield$. Усередині середовища є сферична порожнина. Знайдіть
	у~центрі сфери напруженість $\Efield'$ поля, створеного поляризаційними зарядами,
	індукованими на поверхні сфери, вважаючи, що вектор поляризації $\vect{P}$ всюди (за
	винятком порожнини) має стале значення.
	\begin{solution}
		$\Efield' = \frac{4\pi}{3}\vect P = \frac{\epsilon - 1}{3}\Efield$.
	\end{solution}
\end{problem}

\begin{problem}\label{prb:DielectricSlab}
	Велику плоску пластину з~діелектрика з~проникністю $\epsilon$ вміщено в~однорідне
	зовнішнє поле $\Efield_0$: (а)~перпендикулярно до поля; (б)~паралельно полю. Знайдіть
	$\Efield$, $\Dfield$ і~$\vect P$ усередині пластини та густину зв'язаних зарядів на її
	великих гранях.
	\begin{solution}
		(а) $D = E_0$, $E = E_0/\epsilon$, $P = (\epsilon-1)E_0/(4\pi\epsilon)$, $\sigma' = \pm P$;
		(б) $E = E_0$, $D = \epsilon E_0$, $P = (\epsilon-1)E_0/4\pi$, на великих гранях $\sigma' = 0$.
	\end{solution}
\end{problem}

\begin{problem}\label{prb:BoundAroundQ}
	Точковий заряд $Q$ розташований у~нескінченному однорідному діелектрику з~проникністю
	$\epsilon$. Знайдіть зв'язаний заряд на поверхні малої сфери, яка оточує заряд
	(тобто зв'язаний заряд, що \enquote{прилипає} до заряду $Q$). Чи є в~діелектрику
	об'ємний зв'язаний заряд?
	\begin{solution}
		$q' = -Q(\epsilon-1)/\epsilon$; $\rho' = -\divg\vect P = 0$ при $r > 0$.
	\end{solution}
\end{problem}

\begin{problem}
	Силова лінія поля виходить зі скла ($\epsilon = 6$) у~повітря, утворюючи в~склі кут
	$30^\circ$ з~нормаллю до поверхні. Під яким кутом до нормалі вона йде в~повітрі? Вільних
	зарядів на поверхні немає.
	\begin{solution}
		$\tg\alpha_2 = \tg 30^\circ/6 \approx 0{,}096$, $\alpha_2 \approx 5{,}5^\circ$.
	\end{solution}
\end{problem}

\begin{problem}\label{prb:WaterVapour}
	Оцініть діелектричну проникність водяної пари за температури $100^\circ$C і~тиску
	$1$~атм, якщо дипольний момент молекули \ce{H2O} $p_0 = 1{,}84$~Д, а~її електронна
	поляризовність $\alpha_e \approx 1{,}5\cdot10^{-24}$~см$^3$. Порівняйте орієнтаційний
	та електронний внески. (Виміряне значення $\epsilon \approx 1{,}0059$.)
	\begin{solution}
		$n = p/kT \approx 2\cdot10^{19}$~см$^{-3}$; орієнтаційний внесок $4\pi n p_0^2/3kT \approx 5{,}4\cdot10^{-3}$,
		електронний $4\pi n\alpha_e \approx 3{,}7\cdot10^{-4}$; $\epsilon \approx 1{,}0058$.
	\end{solution}
\end{problem}

\begin{Questions}
	\begin{quizlist}
		\item Чим діелектрик принципово відрізняється від провідника з~погляду наявних у~ньому
		зарядів? Дайте означення вільних та зв'язаних (поляризаційних) зарядів.
		\item Чому мікрополе $\Efield_\text{micro}$ у~речовині неможливо виміряти
		макроскопічними приладами? Дайте означення макроскопічного (усередненого) поля через
		фізично нескінченно малий об'єм $\Delta V$.
		\item Що таке власний дипольний момент молекули і~чому не всі молекули його мають?
		Поясніть на прикладі \ce{H2O} (полярна) і~\ce{CO2} (неполярна), чому симетрія
		молекули визначає наявність дипольного моменту.
		\item Що таке наведений (індукований) дипольний момент неполярної молекули? Запишіть
		формулу $\vect p = \alpha\Efield$. Чому електронна поляризація встигає за полем навіть
		на оптичних частотах?
		\item Чим орієнтаційна поляризація полярних діелектриків відрізняється від електронної
		поляризації неполярних? Чому тепловий рух заважає повній орієнтації диполів уздовж
		поля?
		\item Чому власне поле зв'язаних зарядів $\Efield'$ у~поляризованому діелектрику, як
		правило, послаблює зовнішнє поле? Чому в~пластині, паралельній полю, поле не
		змінюється?
		\item Що таке електрострикція? Чому вона майже не застосовується як самостійне явище
		в~техніці?
		\item Дайте означення вектора поляризації $\vect P$ і~сприйнятливості $\chi$. Які діелектрики називають лінійними, а~які~--- нелінійними (анізотропні кристали, сильні поля, сегнетоелектрики)?

		\item Виведіть формулу $P_n = \sigma'$ для однорідно поляризованого діелектрика,
		розглядаючи косокутний паралелепіпед, вирізаний уздовж вектора поляризації.
		\item Виведіть теорему Гаусса для вектора поляризації в~інтегральній $\oiint P_n\,dS =
			-\iiint \rho'\,dV$ та диференціальній $\divg\vect P = -\rho'$ формах. Чому
		в~однорідно поляризованому діелектрику об'ємний зв'язаний заряд відсутній?
		\item Виведіть означення вектора електричної індукції $\Dfield = \Efield + 4\pi\vect
			P$ з~теореми Гаусса для $\Efield$, що враховує і~вільні, і~зв'язані заряди. Запишіть
		теорему Гаусса для $\Dfield$ в~інтегральній та диференціальній формах.
		\item Чому вектор $\Dfield$ не є силовим полем? У~якому сенсі його джерелами є лише
		вільні заряди і~коли $\Dfield$ можна знайти лише за ними?
		\item Виведіть співвідношення $\Dfield = \epsilon\Efield$ та $\epsilon = 1+4\pi\chi$. Чому поле точкового заряду в~однорідному діелектрику слабкіше, ніж у~вакуумі, рівно в~$\epsilon$ разів, а~поле в~діелектричній кулі в~однорідному полі~--- в~$(\epsilon+2)/3$ раза?

		\item Виведіть граничні умови $D_{2n}-D_{1n}=4\pi\sigma$ і~$E_{1\tau}=E_{2\tau}$ на межі двох діелектриків. Чому $D_n$ без вільних зарядів неперервна, а~$E_n$~--- ні?

		\item Виведіть закон заломлення силових ліній $\tg\alpha_2/\tg\alpha_1 =
			\epsilon_2/\epsilon_1$ на межі двох діелектриків. Чому лінії поля притискаються до
		поверхні середовища з~більшим $\epsilon$?
		\item Які три мікроскопічні механізми поляризовності молекули існують (електронна,
		іонна, орієнтаційна)? Чим вони відрізняються за швидкодією та залежністю від частоти
		й~температури?
		\item Виведіть закон Ланжевена--Дебая $P = np_0^2E/(3kT)$ для орієнтаційної поляризації
		полярного діелектрика, використовуючи розподіл Больцмана для кутів орієнтації
		диполів. За якої умови справедливе лінійне наближення $p_0E\ll kT$?
		\item Виведіть формулу $\epsilon = 1+4\pi np_0^2/(3kT)$ для полярного газу. Яким внеском у~ній знехтувано? Поясніть, чому спадання
		$\epsilon$ з~ростом температури є ознакою полярного діелектрика.
		\item Що таке локальне поле $\Efield_\text{лок}$ і~чим воно відрізняється від
		середнього поля $\Efield$ у~густому середовищі? Виведіть формулу Лоренца
		$\Efield_\text{лок} = \Efield + 4\pi\vect P/3$ для сфери Лоренца.
		\item Виведіть формулу Клаузіуса--Моссотті $(\epsilon-1)/(\epsilon+2) = 4\pi
			n\alpha/3$. У~чому її головна цінність порівняно з~наївним співвідношенням
		$\chi=n\alpha$? При якій умові вона зводиться до останнього?
		\item Що таке електрет і~чому в~ньому $\divg\Dfield\equiv0$ всюди, попри те що поле
		$\Dfield$ поза його поверхнею не дорівнює нулю? Як виготовляють електрети і~де вони
		застосовуються?
		\item Що таке сегнетоелектрик і~доменна структура? Опишіть петлю гістерезису: що таке
		залишкова поляризація $P_r$ та коерцитивне поле $E_c$? Що відбувається при
		температурі Кюрі?
		\item Що таке п'єзо- та піроелектрики, антисегнетоелектрики й~рідкі кристали? Чому всі піроелектрики є п'єзоелектриками, а~не навпаки? На чому ґрунтується робота LCD-дисплеїв?
	\end{quizlist}
\end{Questions}

\Closesolutionfile{answer}
